\section{Overgroups of a two-block stabilizer}
\textbf{Lemma UB.Overgroups (two-block overgroups).} Let \(\Omega=A\sqcup C\) be a finite set, with \(a=|A|\), \(c=|C|\) positive, and put \(H=S_A\times S_C\). If \(a\ne c\), or \(a=c=1\), then \(H\) is maximal in \(S_\Omega\). If \(a=c\geq2\), its only proper strict overgroup is the unordered two-block stabilizer \[
W=(S_A\times S_C)\rtimes C_2,
\] which is maximal and satisfies \([W:H]=2\). In all cases the subgroup interval \([H,S_\Omega]\) is a chain. A positive weighted tree has at most one edge whose two sides have equal total weight.

\emph{Proof.} The case \(a=c=1\) is immediate. The group \(H\) contains the transpositions within each of its two blocks. If an overgroup element does not preserve their unordered partition, its conjugates of these transpositions include one crossing the original blocks. Together with the original within-block transpositions the support graph is connected, so these transpositions generate \(S_\Omega\). This remains valid when one block is a singleton: a permutation moving that singleton produces a cross-block conjugate from the other block. Thus a proper overgroup must preserve the unordered partition. That stabilizer equals \(H\) when the sizes differ and has index two over it when they agree. Its only further overgroup is \(S_\Omega\), by the same argument.

If an edge bisects the total weight, the side of any other edge away from it is a proper positive-weight subset of one of its halves. It has strictly less than half the total weight and cannot be a bisecting side. ∎
